% GENERATED FILE. Edit canonical lessons, then run npm run generate:books.
% canonical-source-hash: fnv1a64:dabb011ecf5a1e6e
% canonical-lessons: UR-C406-lakri, UR-C406-sulgana, UR-C406-orhna, UR-C406-shamiyana, UR-C406-basera

\chapter{Nights Outdoors}
\label{ch:ur-nights-outdoors}
\hlchaptermodality{406}

\begin{chapteropening}
\textbf{By the end of this chapter:} \emph{I can talk about a night in the open: firewood, blankets and shelter.}

Complete the last lesson of chapter 406: I can talk about a night in the open: firewood, blankets and shelter.
\end{chapteropening}

\section[\texorpdfstring{\ur{لکڑی} (lakṛī)}{لکڑی (lakṛī)}]{\ur{لکڑی} (lakṛī) — wood; firewood}

\label{lesson:UR-C406-lakri}



\begin{quote}
Before the new one: say the Urdu for to send for, then the Urdu for a tip (for service).
\end{quote}

\subsection*{You'll want to know: \ur{لکڑی}}
\textbf{\ur{لکڑی}} — \emph{lakṛī} — \textquotedblleft{}wood; firewood\textquotedblright{}.

\begin{quote}
\ur{ہم} \ur{جنگل} \ur{سے} \ur{لکڑی} \ur{لائے۔}
\end{quote}

\begin{quote}
\emph{ham jangal se lakṛī lāe.}
\end{quote}

\begin{quote}
We brought wood back from the forest.
\end{quote}

Feminine.

\begin{cousinweb}
\textbf{\ur{لکڑی}} (\emph{lakṛī}) is inherited: CDIAL lists the Hindi cousin \emph{lakṛī}, \textquotedblleft{}a stick, firewood\textquotedblright{}, under Sanskrit \emph{lakuṭa}, \textquotedblleft{}a club\textquotedblright{} (CDIAL 10875).
\end{cousinweb}

\begin{grammarlens}[title={What you've built}]
One more word for getting around.
\end{grammarlens}

\subsection*{Your turn}
\begin{itemize}
\raggedright
  \item \emph{Say it:} \emph{lakṛī}
  \item \emph{Say it:} \emph{lakṛī}, once more
  \item \emph{Say it:} \emph{ham jangal se lakṛī lāe.} — We brought wood back from the forest.
  \item \emph{Recall:} say the Urdu for to send for, then the Urdu for a tip (for service), then say \emph{lakṛī} again
\end{itemize}

\begin{tcolorbox}[breakable,colback=teal!4,colframe=teal!35!black,title={Before you move on}]
What does \ur{لکڑی} mean? (\textquotedblleft{}Wood; firewood\textquotedblright{}.) Say it once more.
\end{tcolorbox}
\section[\texorpdfstring{\ur{سلگانا} (sulgānā)}{سلگانا (sulgānā)}]{\ur{سلگانا} (sulgānā) — to kindle (a fire); to light}

\label{lesson:UR-C406-sulgana}



\begin{quote}
Before the new one: say the Urdu for a tip (for service), then the Urdu for wood.
\end{quote}

\subsection*{You'll want to know: \ur{سلگانا}}
\textbf{\ur{سلگانا}} — \emph{sulgānā} — \textquotedblleft{}to kindle (a fire); to light\textquotedblright{}.

\begin{quote}
\ur{سارا} \ur{نے} \ur{لکڑی} \ur{سے} \ur{آگ} \ur{سلگائی۔}
\end{quote}

\begin{quote}
\emph{sārā ne lakṛī se āg sulgāī.}
\end{quote}

\begin{quote}
Sara got a fire going with the wood.
\end{quote}

Past, with \ur{نے} (he or she …): \ur{اس} \ur{نے} \ur{چولہا} \ur{سلگایا}, \emph{us ne chūlhā sulgāyā}.

\begin{cousinweb}
\textbf{\ur{سلگانا}} (\emph{sulgānā}) is built on \emph{sulagnā}, \textquotedblleft{}to smoulder, to catch fire\textquotedblright{}, which CDIAL lists under Sanskrit \emph{saṃlagna}, \textquotedblleft{}attached, stuck together\textquotedblright{} (CDIAL 12999): this is the form for setting a fire going.
\end{cousinweb}

\begin{grammarlens}[title={What you've built}]
One more word for getting around.
\end{grammarlens}

\subsection*{Your turn}
\begin{itemize}
\raggedright
  \item \emph{Say it:} \emph{sulgānā}
  \item \emph{Say it:} \emph{sulgānā}, once more
  \item \emph{Say it:} \emph{sārā ne lakṛī se āg sulgāī.} — Sara got a fire going with the wood.
  \item \emph{Recall:} say the Urdu for a tip (for service), then the Urdu for wood, then say \emph{sulgānā} again
\end{itemize}

\begin{tcolorbox}[breakable,colback=teal!4,colframe=teal!35!black,title={Before you move on}]
What does \ur{سلگانا} mean? (\textquotedblleft{}To kindle (a fire); to light\textquotedblright{}.) Say it once more.
\end{tcolorbox}
\section[\texorpdfstring{\ur{اوڑھنا} (oṛhnā)}{اوڑھنا (oṛhnā)}]{\ur{اوڑھنا} (oṛhnā) — to wrap oneself in (a shawl or a blanket)}

\label{lesson:UR-C406-orhna}



\begin{quote}
Before the new one: say the Urdu for wood, then the Urdu for to kindle (a fire).
\end{quote}

\subsection*{You'll want to know: \ur{اوڑھنا}}
\textbf{\ur{اوڑھنا}} — \emph{oṛhnā} — \textquotedblleft{}to wrap oneself in (a shawl or a blanket)\textquotedblright{}.

\begin{quote}
\ur{سردی} \ur{تھی،} \ur{اس} \ur{لیے} \ur{ہم} \ur{نے} \ur{کمبل} \ur{اوڑھے۔}
\end{quote}

\begin{quote}
\emph{sardī thī, is liye ham ne kambal oṛhe.}
\end{quote}

\begin{quote}
It was cold, so we wrapped ourselves in blankets.
\end{quote}

Past, with \ur{نے} (he or she …): \ur{اس} \ur{نے} \ur{چادر} \ur{اوڑھی}, \emph{us ne chādar oṛhī}.

\begin{cousinweb}
\textbf{\ur{اوڑھنا}} (\emph{oṛhnā}) is inherited: CDIAL lists the Hindi cousin \emph{oṛhnā}, \textquotedblleft{}to cover\textquotedblright{}, under \emph{ōḍḍh}, a form it reconstructs (CDIAL 2547).
\end{cousinweb}

\begin{grammarlens}[title={What you've built}]
One more word for getting around.
\end{grammarlens}

\subsection*{Your turn}
\begin{itemize}
\raggedright
  \item \emph{Say it:} \emph{oṛhnā}
  \item \emph{Say it:} \emph{oṛhnā}, once more
  \item \emph{Say it:} \emph{sardī thī, is liye ham ne kambal oṛhe.} — It was cold, so we wrapped ourselves in blankets.
  \item \emph{Recall:} say the Urdu for wood, then the Urdu for to kindle (a fire), then say \emph{oṛhnā} again
\end{itemize}

\begin{tcolorbox}[breakable,colback=teal!4,colframe=teal!35!black,title={Before you move on}]
What does \ur{اوڑھنا} mean? (\textquotedblleft{}To wrap oneself in (a shawl or a blanket)\textquotedblright{}.) Say it once more.
\end{tcolorbox}
\section[\texorpdfstring{\ur{شامیانہ} (shāmiyānā)}{شامیانہ (shāmiyānā)}]{\ur{شامیانہ} (shāmiyānā) — a canopy; a large awning}

\label{lesson:UR-C406-shamiyana}



\begin{quote}
Before the new one: say the Urdu for to kindle (a fire), then the Urdu for to wrap oneself in.
\end{quote}

\subsection*{You'll want to know: \ur{شامیانہ}}
\textbf{\ur{شامیانہ}} — \emph{shāmiyānā} — \textquotedblleft{}a canopy; a large awning\textquotedblright{}.

\begin{quote}
\ur{میدان} \ur{میں} \ur{ایک} \ur{بڑا} \ur{شامیانہ} \ur{تھا۔}
\end{quote}

\begin{quote}
\emph{maidān meṅ ek baṛā shāmiyānā thā.}
\end{quote}

\begin{quote}
There was a big canopy on the field.
\end{quote}

Masculine.

\begin{cousinweb}
\textbf{\ur{شامیانہ}} (\emph{shāmiyānā}) is usually traced to Persian \emph{shāmiyāna}, \textquotedblleft{}a canopy\textquotedblright{}; weddings and camps put one up against the sun and the rain.
\end{cousinweb}

\begin{grammarlens}[title={What you've built}]
One more word for getting around.
\end{grammarlens}

\subsection*{Your turn}
\begin{itemize}
\raggedright
  \item \emph{Say it:} \emph{shāmiyānā}
  \item \emph{Say it:} \emph{shāmiyānā}, once more
  \item \emph{Say it:} \emph{maidān meṅ ek baṛā shāmiyānā thā.} — There was a big canopy on the field.
  \item \emph{Recall:} say the Urdu for to kindle (a fire), then the Urdu for to wrap oneself in, then say \emph{shāmiyānā} again
\end{itemize}

\begin{tcolorbox}[breakable,colback=teal!4,colframe=teal!35!black,title={Before you move on}]
What does \ur{شامیانہ} mean? (\textquotedblleft{}A canopy; a large awning\textquotedblright{}.) Say it once more.
\end{tcolorbox}
\section[\texorpdfstring{\ur{بسیرا} (baserā)}{بسیرا (baserā)}]{\ur{بسیرا} (baserā) — a night's shelter; a roost}

\label{lesson:UR-C406-basera}



\begin{quote}
Before the new one: say the Urdu for to wrap oneself in, then the Urdu for a canopy.
\end{quote}

\subsection*{You'll want to know: \ur{بسیرا}}
\textbf{\ur{بسیرا}} — \emph{baserā} — \textquotedblleft{}a night's shelter; a roost\textquotedblright{}.

\begin{quote}
\ur{پرندے} \ur{شام} \ur{کو} \ur{درختوں} \ur{پر} \ur{بسیرا} \ur{کرتے} \ur{ہیں۔}
\end{quote}

\begin{quote}
\emph{parinde shām ko darakhtoṅ par baserā karte haiṅ.}
\end{quote}

\begin{quote}
In the evening the birds roost in the trees.
\end{quote}

Masculine.

\begin{cousinweb}
\textbf{\ur{بسیرا}} (\emph{baserā}) is inherited: CDIAL lists the Hindi cousin \emph{baserā}, \textquotedblleft{}a place of staying\textquotedblright{}, under \emph{vāsakara}, a form it reconstructs (CDIAL 11594).
\end{cousinweb}

\begin{grammarlens}[title={What you've built}]
One more word for getting around.
\end{grammarlens}

\subsection*{Your turn}
\begin{itemize}
\raggedright
  \item \emph{Say it:} \emph{baserā}
  \item \emph{Say it:} \emph{baserā}, once more
  \item \emph{Say it:} \emph{parinde shām ko darakhtoṅ par baserā karte haiṅ.} — In the evening the birds roost in the trees.
  \item \emph{Recall:} say the Urdu for to wrap oneself in, then the Urdu for a canopy, then say \emph{baserā} again
\end{itemize}

\begin{tcolorbox}[breakable,colback=teal!4,colframe=teal!35!black,title={Before you move on}]
What does \ur{بسیرا} mean? (\textquotedblleft{}A night's shelter; a roost\textquotedblright{}.) Say it once more.
\end{tcolorbox}
